% ============================================================
% fig_ambiguedad.tex: Ambigüedad semántica según contexto
% Uso: % ============================================================
% fig_ambiguedad.tex: Ambigüedad semántica según contexto
% Uso: % ============================================================
% fig_ambiguedad.tex: Ambigüedad semántica según contexto
% Uso: % ============================================================
% fig_ambiguedad.tex: Ambigüedad semántica según contexto
% Uso: \input{figures/fig_ambiguedad}
% ============================================================
\begin{figure}[H]
\centering
\begin{tikzpicture}[
    word/.style={draw=gray!40, rounded corners, fill=gray!8,
                 minimum width=1.1cm, minimum height=0.65cm, font=\small},
    key/.style={draw=blue!60, rounded corners, fill=blue!12,
                minimum width=1.1cm, minimum height=0.65cm, font=\small\bfseries},
    meaning/.style={draw, rounded corners, fill=white,
                    minimum width=2.6cm, minimum height=0.7cm,
                    font=\footnotesize, align=center},
    arrow/.style={-{Stealth}, thick},
    ctx/.style={font=\footnotesize\itshape, gray}
]

% ── Oración 1 ────────────────────────────────────────────────
\node[ctx] at (3.5, 5.6) {Contexto A:};
\foreach \w/\i in {Fui/0, al/1, banco/2, a/3, depositar/4, dinero/5}{
    \pgfmathparse{int(\i==2)}
    \ifnum\pgfmathresult=1
        \node[key]  at (\i*1.18, 4.95) {\w};
    \else
        \node[word] at (\i*1.18, 4.95) {\w};
    \fi
}
\draw[arrow, blue!60] (2.36, 4.62) -- (2.36, 4.0);
\node[meaning, fill=green!10, draw=green!50!black]
    at (2.36, 3.6) {institución\\financiera};

% ── Oración 2 ────────────────────────────────────────────────
\node[ctx] at (3.5, 2.9) {Contexto B:};
\foreach \w/\i in {Me/0, senté/1, en/2, el/3, banco/4, del/5, parque/6}{
    \pgfmathparse{int(\i==4)}
    \ifnum\pgfmathresult=1
        \node[key]  at (\i*1.1, 2.3) {\w};
    \else
        \node[word] at (\i*1.1, 2.3) {\w};
    \fi
}
\draw[arrow, blue!60] (4.4, 1.97) -- (4.4, 1.35);
\node[meaning, fill=orange!10, draw=orange!60]
    at (4.4, 0.95) {asiento\\exterior};

% ── Anotación central ────────────────────────────────────────
\node[draw=blue!20, fill=blue!5, rounded corners, font=\footnotesize,
      inner sep=6pt, align=center] at (7.8, 2.8)
    {\textbf{``banco''} es la misma\\palabra en ambos casos.\\El significado\\
     solo emerge del\\contexto circundante.};

\draw[dashed, gray!50] (5.9, 3.6) -- (6.8, 3.2);
\draw[dashed, gray!50] (5.8, 0.95) -- (6.8, 2.4);

\end{tikzpicture}
\caption{Ambigüedad semántica léxica. La palabra \textit{banco} tiene
significados completamente distintos según el contexto oracional. Un modelo
que procese cada token de forma independiente, sin considerar las palabras
circundantes, no puede resolver esta ambigüedad.}
\label{fig:ambiguedad}
\end{figure}

% ============================================================
\begin{figure}[H]
\centering
\begin{tikzpicture}[
    word/.style={draw=gray!40, rounded corners, fill=gray!8,
                 minimum width=1.1cm, minimum height=0.65cm, font=\small},
    key/.style={draw=blue!60, rounded corners, fill=blue!12,
                minimum width=1.1cm, minimum height=0.65cm, font=\small\bfseries},
    meaning/.style={draw, rounded corners, fill=white,
                    minimum width=2.6cm, minimum height=0.7cm,
                    font=\footnotesize, align=center},
    arrow/.style={-{Stealth}, thick},
    ctx/.style={font=\footnotesize\itshape, gray}
]

% ── Oración 1 ────────────────────────────────────────────────
\node[ctx] at (3.5, 5.6) {Contexto A:};
\foreach \w/\i in {Fui/0, al/1, banco/2, a/3, depositar/4, dinero/5}{
    \pgfmathparse{int(\i==2)}
    \ifnum\pgfmathresult=1
        \node[key]  at (\i*1.18, 4.95) {\w};
    \else
        \node[word] at (\i*1.18, 4.95) {\w};
    \fi
}
\draw[arrow, blue!60] (2.36, 4.62) -- (2.36, 4.0);
\node[meaning, fill=green!10, draw=green!50!black]
    at (2.36, 3.6) {institución\\financiera};

% ── Oración 2 ────────────────────────────────────────────────
\node[ctx] at (3.5, 2.9) {Contexto B:};
\foreach \w/\i in {Me/0, senté/1, en/2, el/3, banco/4, del/5, parque/6}{
    \pgfmathparse{int(\i==4)}
    \ifnum\pgfmathresult=1
        \node[key]  at (\i*1.1, 2.3) {\w};
    \else
        \node[word] at (\i*1.1, 2.3) {\w};
    \fi
}
\draw[arrow, blue!60] (4.4, 1.97) -- (4.4, 1.35);
\node[meaning, fill=orange!10, draw=orange!60]
    at (4.4, 0.95) {asiento\\exterior};

% ── Anotación central ────────────────────────────────────────
\node[draw=blue!20, fill=blue!5, rounded corners, font=\footnotesize,
      inner sep=6pt, align=center] at (7.8, 2.8)
    {\textbf{``banco''} es la misma\\palabra en ambos casos.\\El significado\\
     solo emerge del\\contexto circundante.};

\draw[dashed, gray!50] (5.9, 3.6) -- (6.8, 3.2);
\draw[dashed, gray!50] (5.8, 0.95) -- (6.8, 2.4);

\end{tikzpicture}
\caption{Ambigüedad semántica léxica. La palabra \textit{banco} tiene
significados completamente distintos según el contexto oracional. Un modelo
que procese cada token de forma independiente, sin considerar las palabras
circundantes, no puede resolver esta ambigüedad.}
\label{fig:ambiguedad}
\end{figure}

% ============================================================
\begin{figure}[H]
\centering
\begin{tikzpicture}[
    word/.style={draw=gray!40, rounded corners, fill=gray!8,
                 minimum width=1.1cm, minimum height=0.65cm, font=\small},
    key/.style={draw=blue!60, rounded corners, fill=blue!12,
                minimum width=1.1cm, minimum height=0.65cm, font=\small\bfseries},
    meaning/.style={draw, rounded corners, fill=white,
                    minimum width=2.6cm, minimum height=0.7cm,
                    font=\footnotesize, align=center},
    arrow/.style={-{Stealth}, thick},
    ctx/.style={font=\footnotesize\itshape, gray}
]

% ── Oración 1 ────────────────────────────────────────────────
\node[ctx] at (3.5, 5.6) {Contexto A:};
\foreach \w/\i in {Fui/0, al/1, banco/2, a/3, depositar/4, dinero/5}{
    \pgfmathparse{int(\i==2)}
    \ifnum\pgfmathresult=1
        \node[key]  at (\i*1.18, 4.95) {\w};
    \else
        \node[word] at (\i*1.18, 4.95) {\w};
    \fi
}
\draw[arrow, blue!60] (2.36, 4.62) -- (2.36, 4.0);
\node[meaning, fill=green!10, draw=green!50!black]
    at (2.36, 3.6) {institución\\financiera};

% ── Oración 2 ────────────────────────────────────────────────
\node[ctx] at (3.5, 2.9) {Contexto B:};
\foreach \w/\i in {Me/0, senté/1, en/2, el/3, banco/4, del/5, parque/6}{
    \pgfmathparse{int(\i==4)}
    \ifnum\pgfmathresult=1
        \node[key]  at (\i*1.1, 2.3) {\w};
    \else
        \node[word] at (\i*1.1, 2.3) {\w};
    \fi
}
\draw[arrow, blue!60] (4.4, 1.97) -- (4.4, 1.35);
\node[meaning, fill=orange!10, draw=orange!60]
    at (4.4, 0.95) {asiento\\exterior};

% ── Anotación central ────────────────────────────────────────
\node[draw=blue!20, fill=blue!5, rounded corners, font=\footnotesize,
      inner sep=6pt, align=center] at (7.8, 2.8)
    {\textbf{``banco''} es la misma\\palabra en ambos casos.\\El significado\\
     solo emerge del\\contexto circundante.};

\draw[dashed, gray!50] (5.9, 3.6) -- (6.8, 3.2);
\draw[dashed, gray!50] (5.8, 0.95) -- (6.8, 2.4);

\end{tikzpicture}
\caption{Ambigüedad semántica léxica. La palabra \textit{banco} tiene
significados completamente distintos según el contexto oracional. Un modelo
que procese cada token de forma independiente, sin considerar las palabras
circundantes, no puede resolver esta ambigüedad.}
\label{fig:ambiguedad}
\end{figure}

% ============================================================
\begin{figure}[H]
\centering
\begin{tikzpicture}[
    word/.style={draw=gray!40, rounded corners, fill=gray!8,
                 minimum width=1.1cm, minimum height=0.65cm, font=\small},
    key/.style={draw=blue!60, rounded corners, fill=blue!12,
                minimum width=1.1cm, minimum height=0.65cm, font=\small\bfseries},
    meaning/.style={draw, rounded corners, fill=white,
                    minimum width=2.6cm, minimum height=0.7cm,
                    font=\footnotesize, align=center},
    arrow/.style={-{Stealth}, thick},
    ctx/.style={font=\footnotesize\itshape, gray}
]

% ── Oración 1 ────────────────────────────────────────────────
\node[ctx] at (3.5, 5.6) {Contexto A:};
\foreach \w/\i in {Fui/0, al/1, banco/2, a/3, depositar/4, dinero/5}{
    \pgfmathparse{int(\i==2)}
    \ifnum\pgfmathresult=1
        \node[key]  at (\i*1.18, 4.95) {\w};
    \else
        \node[word] at (\i*1.18, 4.95) {\w};
    \fi
}
\draw[arrow, blue!60] (2.36, 4.62) -- (2.36, 4.0);
\node[meaning, fill=green!10, draw=green!50!black]
    at (2.36, 3.6) {institución\\financiera};

% ── Oración 2 ────────────────────────────────────────────────
\node[ctx] at (3.5, 2.9) {Contexto B:};
\foreach \w/\i in {Me/0, senté/1, en/2, el/3, banco/4, del/5, parque/6}{
    \pgfmathparse{int(\i==4)}
    \ifnum\pgfmathresult=1
        \node[key]  at (\i*1.1, 2.3) {\w};
    \else
        \node[word] at (\i*1.1, 2.3) {\w};
    \fi
}
\draw[arrow, blue!60] (4.4, 1.97) -- (4.4, 1.35);
\node[meaning, fill=orange!10, draw=orange!60]
    at (4.4, 0.95) {asiento\\exterior};

% ── Anotación central ────────────────────────────────────────
\node[draw=blue!20, fill=blue!5, rounded corners, font=\footnotesize,
      inner sep=6pt, align=center] at (7.8, 2.8)
    {\textbf{``banco''} es la misma\\palabra en ambos casos.\\El significado\\
     solo emerge del\\contexto circundante.};

\draw[dashed, gray!50] (5.9, 3.6) -- (6.8, 3.2);
\draw[dashed, gray!50] (5.8, 0.95) -- (6.8, 2.4);

\end{tikzpicture}
\caption{Ambigüedad semántica léxica. La palabra \textit{banco} tiene
significados completamente distintos según el contexto oracional. Un modelo
que procese cada token de forma independiente, sin considerar las palabras
circundantes, no puede resolver esta ambigüedad.}
\label{fig:ambiguedad}
\end{figure}
